\section{Model Description}
\label{sec:model}

\subsection{Field Definitions}

The DBS micro-continuum formulation \citep{brinkman1949calculation,soulaine2016micro,
soulaine2024microcontinuum} represents the pore medium as a continuum in which
a cell-centred porosity field $\epsilon$ varies smoothly from
zero (solid) to unity (open pore). The steady incompressible flow is governed
by the Darcy-Brinkman equations:
\begin{align}
\nabla \cdot \bm{u}_D &= 0,
\label{eq:continuity} \\
-\nabla p + \nabla \cdot \bigl(\mu_{\mathrm{eff}} \nabla \bm{u}\bigr)
- \frac{\mu}{K(\epsilon)}\,\bm{u}_D &= \bm{0},
\label{eq:brinkman}
\end{align}
where $\bm{u}_D = \epsilon\bm{u}$ is the Darcy velocity,
$\mu$ is the fluid viscosity, $\mu_{\mathrm{eff}}$ is an effective
viscosity that smoothly recovers the Stokes equation in open-pore cells
\citep{soulaine2016micro}, and $K(\epsilon)$ is a permeability
function. The drag coefficient $\mu/K(\epsilon)$ therefore acts
as a large resistance term in solid-like cells and vanishes in open-pore cells,
so Eq.~\eqref{eq:brinkman} reduces to the Stokes equation.

For an image-derived cell, let $f_k$ be the geometric fraction of source voxels
assigned to segmented phase $k$, with $\sum_k f_k=1$, and let $\phi_k$ be the
end-member porosity assigned to that phase. The cell DBS porosity is
\begin{equation}
\epsilon=\sum_k f_k\phi_k.
\label{eq:phase_porosity}
\end{equation}
The fraction $f_k$ describes phase occupancy, whereas $\phi_k$ specifies the
water-filled fraction of that phase. A cell wholly occupied by a microporous
phase therefore has $f_{\mathrm{micro}}=1$ and
$\epsilon=\phi_{\mathrm{micro}}<1$. Nominal solid is assigned a small non-zero
porosity as a numerical DBS floor.

The mobile-water fraction $\theta$ is the mobile-water volume per bulk cell
volume. It governs aqueous storage, diffusion, and PHREEQC solution volume,
whereas $\epsilon$ governs DBS flow resistance and relates intrinsic velocity
to Darcy velocity. The saturated calculations reported here set
$\theta=\epsilon$ in chemically active cells. Cells with
$\theta\leq\theta_{\mathrm{mask}}$ are excluded from chemistry.

Reactive surface area is specified independently as the cell-centred bulk-volume
density $a_{\mathrm{surf}}$ (m$^2$\,m$^{-3}_{\mathrm{bulk}}$). All aqueous
species, $\theta$, $a_{\mathrm{surf}}$, and adsorbed surface fields are
cell-centred; transport fluxes are evaluated on cell faces.

\subsection{Reactive Surface Area Field}
\label{sec:aSurf}

Reactive area is present only for selected interfaces and phases assigned a
non-zero intrinsic density. Its resolved and unresolved contributions are
carried in the fixed cell-centred field $a_{\mathrm{surf}}$ (units:
m$^2\,\mathrm{m}^{-3}_{\mathrm{bulk}}$), the sum of two terms:
\begin{equation}
a_{\mathrm{surf}} = a_{\mathrm{surf}}^{\mathrm{interface}} +
a_{\mathrm{surf}}^{\mathrm{intrinsic}}.
\label{eq:asurf}
\end{equation}

The interface term is reconstructed for an ordered pair of segmented labels.
One label is the owner and the second is the other phase; their union restricts
the extracted interface, and the owner identifies the side used for area
assignment. In \texttt{occupancyGaussian}, the Gaussian-filtered owner
indicator is normalised by the Gaussian-filtered pair mask, and the $0.5$
level is extracted.

Three-dimensional triangle area is divided by bulk cell volume to obtain
$a_{\mathrm{surf}}^{\mathrm{interface}}$. In quasi-two-dimensional images,
contour length is multiplied by physical out-of-plane thickness before
division by cell volume. For quasi-two-dimensional cell-field construction,
\texttt{occupancyGaussian} uses weighted redistribution over owner-side
first-shell contact cells. Pieces with empty support contribute zero area
to the deposited field. Detailed support and partition rules are given in
Online Resource~1.

\begin{table*}[htbp]
    \centering
    \caption{Interface reconstruction parameters. All cases use
    \texttt{occupancyGaussian} with an isovalue of $0.5$.
    The ordered pair is owner $\rightarrow$ other; the owner defines the
    indicator used for reconstruction. In cell-field constructions,
    reconstructed area is assigned on the owner side.}
    \label{tab:surfaceConstruction}
    \scriptsize
    \setlength{\tabcolsep}{3pt}
    \begin{tabular}{@{}p{0.28\textwidth}p{0.12\textwidth}p{0.48\textwidth}@{}}
        \toprule
        Case & $\sigma$ (vox) & Pair (owner $\rightarrow$ other) \\
        \midrule
        Circular post & 1.50
            & exterior/pore 255 $\rightarrow$ post/solid 0 \\
        Sphere & 0.75
            & exterior/pore 255 $\rightarrow$ sphere/solid 0 \\
        HM12 solid-grain DBS & 0.75
            & pore 255 $\rightarrow$ solid 0 \\
        2D four-fracture & 0.75
            & micropore 1 $\rightarrow$ pore 2 \\
        3D Kelvin & 0.75
            & micropore 2 $\rightarrow$ pore 1 \\
        3D Estaillades & 0.74
            & pore 1 $\rightarrow$ solid 3 \\
        3D external area & 0.74
            & pore 1 $\rightarrow$ remapped solid 3 \\
        \bottomrule
    \end{tabular}
\end{table*}

The interface term captures only boundaries between grain and fluid that are
visible at the resolved image scale. It does not detect
reactive area that lies entirely \emph{within} a voxel, such as sub-voxel
surface roughness of a grain or the internal pore network of a microporous
phase. Those surfaces are below the image resolution by construction.

The second term, $a_{\mathrm{surf}}^{\mathrm{intrinsic}}$, accounts for this
sub-voxel contribution. An end-member surface density $s_k$
(units: m$^2$\,m$^{-3}_{k}$) is prescribed per unit volume of each segmented
phase $k$ (for example pore, micropore, or dense solid). It may represent
sub-voxel surface such as the
Brunauer-Emmett-Teller (BET) surface area of a microporous grain measured by
N$_2$ adsorption, or the surface roughness of a dense solid grain measured on
powdered material. Mass-specific measurements are converted to phase-volume
densities using the corresponding mass per unit phase volume.

The intrinsic contribution is assembled by weighting each phase density by
its geometric fraction $f_k$ within the computational cell:
\begin{equation}
a_{\mathrm{surf}}^{\mathrm{intrinsic}} = \sum_k f_k\,s_k.
\label{eq:asurf_intrinsic}
\end{equation}
Thus $f_k$ converts the phase-volume density $s_k$ to a bulk-volume cell
contribution. A density already expressed per bulk cell volume represents
that contribution directly. The segmented image supplies $f_k$, whereas
$s_k$ is prescribed from complementary measurements or finer-scale
characterisation.

$s_k$ may be non-zero for any label, including labels selected for interface
reconstruction. The same segmented label may contribute through both terms:
$a_{\mathrm{surf}}^{\mathrm{interface}}$ gives the image-scale outer surface
of the grain, and $a_{\mathrm{surf}}^{\mathrm{intrinsic}}$ adds the
sub-voxel surface area that the image resolution cannot capture. If
$s_k=0$ for all labels, the total reduces to the purely geometric image-scale
surface. Both contributions are pre-computed using the segmented image and
prescribed phase properties before the flow simulation begins and are held
constant throughout the reactive-transport calculation.

Both contributions enter the geochemical module through the same bulk-volume
field. The \texttt{phreeqcMixture} builds used for the reported calculations
supply surface area per representative bulk litre to PHREEQC as
\begin{equation}
a_{\mathrm{surf}}^{\mathrm{bulk,L}} = \frac{a_{\mathrm{surf}}}{10^{3}},
\label{eq:asurf_bulk_litre}
\end{equation}
For surface master $k$, the site density $N_k$ (mol\,m$^{-2}$) gives a
surface-master inventory per bulk volume of $N_k a_{\mathrm{surf}}$
(mol\,m$^{-3}_{\mathrm{bulk}}$). These conversions use the bulk-volume basis
without division by $\theta$.

Surface areas estimated from nano-computed tomography
(nano-CT), micro-computed tomography (micro-CT), focused ion beam scanning
electron microscopy (FIB-SEM), or other image data are first expressed on a
bulk-volume basis (m$^2$\,m$^{-3}_{\mathrm{bulk}}$), consistent with
Eqs.~\eqref{eq:asurf} and~\eqref{eq:asurf_intrinsic}. Image-derived surface areas
reported per microporous-phase volume or per grain mass are converted to this
bulk-volume basis before Eq.~\eqref{eq:asurf_bulk_litre} is applied.

The same fixed porosity field is used for DBS transport and for the PHREEQC
solution volume. Cells in
which $\theta$ falls below the chemistry mask threshold
$\theta_{\mathrm{mask}}$ are treated as chemically inactive and their surface
terms are suppressed.

\subsection{Conservative Aqueous Transport}

For each aqueous species with concentration $C_i$, the conservative DBS
transport equation is written as
\begin{equation}
\frac{\partial (\theta C_i)}{\partial t}
+ \nabla \cdot (\bm{u}_D C_i)
- \nabla \cdot (\theta D_i \nabla C_i) = 0,
\end{equation}
where $\bm{u}_D$ is the Darcy velocity satisfying
Eqs.~\eqref{eq:continuity} and~\eqref{eq:brinkman}, and $D_i$ denotes the
diffusivity used in the transport equation. In fully resolved macro-pore cells
($\theta \approx 1$), $D_i$ is the free molecular diffusivity. In cells that
represent microporous material ($0 < \theta < 1$), $D_i$ is interpreted as an
effective diffusivity, commonly $D_i = D_i^{\mathrm{mol}}/\tau$, where
$D_i^{\mathrm{mol}}$ is the molecular diffusivity and $\tau \geq 1$ is the
tortuosity factor of the micropore network. In either case, the $\theta$ factor
multiplying $D_i$ in the diffusion term accounts for the reduced water volume
available for transport. Reactions enter through the operator-split PHREEQC
step described below, not as a source term in this equation.

For a vector of aqueous species $\bm{C}$, the same balance is
\begin{equation}
\frac{\partial (\theta \bm{C})}{\partial t}
+ \nabla \cdot (\bm{u}_D \bm{C})
- \nabla \cdot (\theta \bm{D} \nabla \bm{C}) = \bm{0},
\end{equation}
where $\bm{D}$ collects the species diffusivities. This notation is used for
the split operators below.

\subsection{Cell-Based Surface Inventory}

For the cell-based unresolved surface carried by the DBS formulation, the
conservative local inventory is expressed
in the current solver field variables as
\begin{equation}
M = \theta C + a_{\mathrm{surf}}\Gamma,
\label{eq:inventory}
\end{equation}
where $\Gamma$ is the cell-centred surface-species concentration
(mol\,m$^{-2}$ of reactive surface). The first term is the aqueous amount per
bulk volume, the second the adsorbed amount per bulk volume. Porosity enters
the aqueous term once, and the bulk-based surface-area field carries no
porosity scaling.

For multiple surface species the local bulk inventory is the sum of the aqueous
and adsorbed contributions over species; single-species notation is used
below.

\subsection{Local Chemistry Coupling}

The chemical update is local to each cell. Over one cycle of transport and
chemistry, the transport step first advances the aqueous species in
conservative finite-volume form. The updated cell state, together with the mobile-water
fraction and cell-based surface state, is then mapped into PHREEQC
\citep{parkhurst2013phreeqc} via the PhreeqcRM interface
\citep{parkhurst2015phreeqcrm}. PHREEQC solves the local chemistry at
\emph{thermodynamic equilibrium}, treating surface complexation as
instantaneous relative to transport within each timestep. This holds when the
surface equilibration time is short compared with the time for advection and
diffusion across a cell. The same interface supports kinetic surface reactions,
which are not used here.

The solver uses first-order Lie splitting \citep{carrayrou2004operator}. The
splitting introduces a global error of $\mathcal{O}(\Delta t)$ in the
operator-split solution; the local error over one step is
$\mathcal{O}(\Delta t^2)$. For the timesteps used here, the advective
Courant-Friedrichs-Lewy (CFL) number remains $O(1)$, so the
operator-splitting error is of the same formal order as the first-order
transport discretisation. The resulting algorithm keeps the conservative
transport update and the local PHREEQC equilibrium update as distinct
operators within each timestep.

At the algorithmic level, one timestep of size $\Delta t$ is:
\begin{align}
\widetilde{\bm{C}}^{\,n+1} &= \mathcal{T}_{\Delta t}(\bm{C}^{n}, \theta, \bm{u}_D, \bm{D}),\\
(\bm{C}^{n+1}, \bm{\Gamma}^{n+1}) &=
\mathcal{R}_{\Delta t}(\widetilde{\bm{C}}^{\,n+1}, \bm{\Gamma}^{n}, \theta, a_{\mathrm{surf}}),
\end{align}
where $\mathcal{T}_{\Delta t}$ is the conservative finite-volume transport
operator (advancing $\theta\bm{C}$ in DBS form) and $\mathcal{R}_{\Delta t}$
is the local PHREEQC equilibrium operator that updates both aqueous and
surface states simultaneously within each cell.

\subsection{Analytical Dilute Limit}

In the dilute linear-equilibrium limit, surface complexation reduces to a
Henry-type isotherm. The standard Henry coefficient is $K_s = \Gamma / C$
(units: m), the ratio of surface concentration $\Gamma$ (mol\,m$^{-2}$) to
aqueous concentration $C$ (mol\,m$^{-3}$):
\begin{equation}
\Gamma = K_s C.
\end{equation}
$K_s$ is obtained from a pair of batch equilibrations at two aqueous
concentrations, $K_s = \Delta\Gamma / \Delta C$, which returns the coefficient
in the field units used by the retardation formula.

Substituting into the inventory Eq.~\eqref{eq:inventory} and taking the
time derivative yields the retarded 1D transport equation
\begin{equation}
\left(
\theta + a_{\mathrm{surf}}K_s
\right)\frac{\partial C}{\partial t}
+ \frac{\partial (u_D C)}{\partial x} = 0.
\end{equation}
The corresponding retardation factor is
\begin{equation}
R = 1 + \frac{a_{\mathrm{surf}}\,K_s}{\theta}.
\label{eq:retardation}
\end{equation}
Both $a_{\mathrm{surf}}$ and $\theta$ can vary spatially, so
Eq.~\eqref{eq:retardation} also gives the local retardation at any point in
a heterogeneous domain. The classical semi-infinite advection--diffusion
solution follows \citet{ogata1961solution}. For comparison with the finite 1D
calculation, the reference problem is instead evaluated on $[0,L]$ with the
same fixed-concentration inlet and zero-gradient outlet as the numerical case,
using the effective velocity $u_D/(\theta R)$ and diffusivity $D/R$.
